\documentclass{amsart}
% For dealing with references we use the comment environment
\usepackage{verbatim}
\newenvironment{reference}{\comment}{\endcomment}
%\newenvironment{reference}{}{}
\newenvironment{slogan}{\comment}{\endcomment}
\newenvironment{history}{\comment}{\endcomment}

% For commutative diagrams we use Xy-pic
\usepackage[all]{xy}

% We use 2cell for 2-commutative diagrams.
\xyoption{2cell}
\UseAllTwocells

% We use multicol for the list of chapters between chapters
\usepackage{multicol}

% This is generally recommended for better output
\usepackage{lmodern}
\usepackage[T1]{fontenc}

% For cross-file-references

% Package for hypertext links:
\usepackage{hyperref}

% For any local file, say "hello.tex" you want to link to please

% Theorem environments.
%
\theoremstyle{plain}
\newtheorem{theorem}[subsection]{Theorem}
\newtheorem{proposition}[subsection]{Proposition}
\newtheorem{lemma}[subsection]{Lemma}

\theoremstyle{definition}
\newtheorem{definition}[subsection]{Definition}
\newtheorem{example}[subsection]{Example}
\newtheorem{exercise}[subsection]{Exercise}
\newtheorem{situation}[subsection]{Situation}

\theoremstyle{remark}
\newtheorem{remark}[subsection]{Remark}
\newtheorem{remarks}[subsection]{Remarks}

\numberwithin{equation}{subsection}

% Macros
%
\def\lim{\mathop{\mathrm{lim}}\nolimits}
\def\colim{\mathop{\mathrm{colim}}\nolimits}
\def\Spec{\mathop{\mathrm{Spec}}}
\def\Hom{\mathop{\mathrm{Hom}}\nolimits}
\def\Ext{\mathop{\mathrm{Ext}}\nolimits}
\def\SheafHom{\mathop{\mathcal{H}\!\mathit{om}}\nolimits}
\def\SheafExt{\mathop{\mathcal{E}\!\mathit{xt}}\nolimits}
\def\Sch{\mathit{Sch}}
\def\Mor{\mathop{\mathrm{Mor}}\nolimits}
\def\Ob{\mathop{\mathrm{Ob}}\nolimits}
\def\Sh{\mathop{\mathit{Sh}}\nolimits}
\def\NL{\mathop{N\!L}\nolimits}
\def\CH{\mathop{\mathrm{CH}}\nolimits}
\def\proetale{{pro\text{-}\acute{e}tale}}
\def\etale{{\acute{e}tale}}
\def\QCoh{\mathit{QCoh}}
\def\Ker{\mathop{\mathrm{Ker}}}
\def\Im{\mathop{\mathrm{Im}}}
\def\Coker{\mathop{\mathrm{Coker}}}
\def\Coim{\mathop{\mathrm{Coim}}}

% Boxtimes
%
\DeclareMathSymbol{\boxtimes}{\mathbin}{AMSa}{"02}

%
% Macros for moduli stacks/spaces
%
\def\QCohstack{\mathcal{QC}\!\mathit{oh}}
\def\Cohstack{\mathcal{C}\!\mathit{oh}}
\def\Spacesstack{\mathcal{S}\!\mathit{paces}}
\def\Quotfunctor{\mathrm{Quot}}
\def\Hilbfunctor{\mathrm{Hilb}}
\def\Curvesstack{\mathcal{C}\!\mathit{urves}}
\def\Polarizedstack{\mathcal{P}\!\mathit{olarized}}
\def\Complexesstack{\mathcal{C}\!\mathit{omplexes}}
% \Pic is the operator that assigns to X its picard group, usage \Pic(X)
% \Picardstack_{X/B} denotes the Picard stack of X over B
% \Picardfunctor_{X/B} denotes the Picard functor of X over B
\def\Pic{\mathop{\mathrm{Pic}}\nolimits}
\def\Picardstack{\mathcal{P}\!\mathit{ic}}
\def\Picardfunctor{\mathrm{Pic}}
\def\Deformationcategory{\mathcal{D}\!\mathit{ef}}

\setlength{\emergencystretch}{3em}
\begin{document}
\title[Stacks Project, Tag 0794]{318 --- Existence of unbounded left derived functors under exact sequential colimits}
\maketitle
\noindent Copyright (C) 2005--2025 Johan de Jong. Stacks Project authors. Source: \href{https://stacks.math.columbia.edu/tag/0794}{Tag 0794}.

\noindent Editorial difficulty note (not author text). Difficulty H2: The proof builds a compatible sequence of bounded-above resolutions whose colimit resolves an arbitrary complex. It handles the essential difficulty that a morphism of colimit resolutions need not come from levelwise maps, using common bounded resolutions and compatible cohomology comparisons. Exactness and preservation of sequential colimits then yield the desired quasi-isomorphism..

\noindent Editorial provenance note (not author text). History: excerpt from revision \texttt{a0444\allowbreak 6e57e\allowbreak c1fbc\allowbreak 252a8\allowbreak 71afc\allowbreak ec775\allowbreak 2fb28\allowbreak 07b14}, derived.tex, lines 9332--9442. Reference presentation was adapted; original-author context and core support blocks were retained or added. The mathematical wording is unchanged. Licensed under GNU FDL 1.2 or later; no invariant sections or cover texts. See ../licenses/COPYING. Context blocks reproduce author definitions and supporting results; they are not additional counted problems.

\begin{situation}
\label{situation-derived-functor}
Here $F : \mathcal{D} \to \mathcal{D}'$ is an exact functor of triangulated
categories and $S$ is a saturated multiplicative
system in $\mathcal{D}$ compatible with the structure
of triangulated category on $\mathcal{D}$.
\end{situation}

\begin{definition}
\label{definition-right-derived-functor-defined}
Assumptions and notation as in
Situation \ref{situation-derived-functor}.
Let $X \in \Ob(\mathcal{D})$.
\begin{enumerate}
\item we say the {\it right derived functor $RF$ is defined at}
$X$ if the ind-object
$$
(X/S) \longrightarrow \mathcal{D}', \quad
(s : X \to X') \longmapsto F(X')
$$
is essentially constant\footnote{For a discussion of when an ind-object
or pro-object of a category is essentially constant we refer to
Categories, Section \href{https://stacks.math.columbia.edu/tag/05PT}{section essentially constant (Tag 05PT)}.};
in this case the value
$Y$ in $\mathcal{D}'$ is called the {\it value of $RF$ at $X$}.
\item we say the {\it left derived functor $LF$ is defined at} $X$
if the pro-object
$$
(S/X) \longrightarrow \mathcal{D}', \quad
(s: X' \to X) \longmapsto F(X')
$$
is essentially constant; in this case the value $Y$ in $\mathcal{D}'$
is called the {\it value of $LF$ at $X$}.
\end{enumerate}
By abuse of notation we often denote the values simply
$RF(X)$ or $LF(X)$.
\end{definition}

\begin{situation}
\label{situation-classical}
Here $F : \mathcal{A} \to \mathcal{B}$ is an additive functor between
abelian categories. This induces exact functors
$$
F : K(\mathcal{A}) \to K(\mathcal{B}), \quad
K^{+}(\mathcal{A}) \to K^{+}(\mathcal{B}), \quad
K^{-}(\mathcal{A}) \to K^{-}(\mathcal{B}).
$$
See Lemma \href{https://stacks.math.columbia.edu/tag/014X}{lemma additive exact homotopy category (Tag 014X)}.
We also denote $F$ the composition $K(\mathcal{A}) \to D(\mathcal{B})$,
$K^{+}(\mathcal{A}) \to D^{+}(\mathcal{B})$, and
$K^{-}(\mathcal{A}) \to D^-(\mathcal{B})$ of $F$ with the localization
functor $K(\mathcal{B}) \to D(\mathcal{B})$, etc. This situation leads
to four derived functors we will consider in the following.
\begin{enumerate}
\item The right derived functor of
$F : K(\mathcal{A}) \to D(\mathcal{B})$
relative to the multiplicative system $\text{Qis}(\mathcal{A})$.
\item The right derived functor of
$F : K^{+}(\mathcal{A}) \to D^{+}(\mathcal{B})$
relative to the multiplicative system $\text{Qis}^{+}(\mathcal{A})$.
\item The left derived functor of
$F : K(\mathcal{A}) \to D(\mathcal{B})$
relative to the multiplicative system $\text{Qis}(\mathcal{A})$.
\item The left derived functor of
$F : K^{-}(\mathcal{A}) \to D^{-}(\mathcal{B})$
relative to the multiplicative system $\text{Qis}^-(\mathcal{A})$.
\end{enumerate}
Each of these cases is an example of
Situation \ref{situation-derived-functor}.
\end{situation}

\begin{proposition}
\label{proposition-left-derived-exists}
Let $F : \mathcal{A} \to \mathcal{B}$ be a right exact functor
of abelian categories. Let $\mathcal{P} \subset \Ob(\mathcal{A})$ be a
subset. Assume
\begin{enumerate}
\item $\mathcal{P}$ contains $0$, is closed under (finite) direct sums,
and every object of $\mathcal{A}$ is a quotient of an
element of $\mathcal{P}$,
\item for any bounded above acyclic complex $P^\bullet$ of
$\mathcal{A}$ with $P^n \in \mathcal{P}$ for all $n$ the
complex $F(P^\bullet)$ is exact,
\item $\mathcal{A}$ and $\mathcal{B}$ have colimits
of systems over $\mathbf{N}$,
\item colimits over $\mathbf{N}$ are exact in both
$\mathcal{A}$ and $\mathcal{B}$, and
\item $F$ commutes with colimits over $\mathbf{N}$.
\end{enumerate}
Then $LF$ is defined on all of $D(\mathcal{A})$.
\end{proposition}

\begin{proof}
By (1) and Lemma \href{https://stacks.math.columbia.edu/tag/05T7}{lemma subcategory left resolution (Tag 05T7)} for any bounded
above complex $K^\bullet$ there exists a quasi-isomorphism
$P^\bullet \to K^\bullet$ with $P^\bullet$ bounded above and
$P^n \in \mathcal{P}$ for all $n$. Suppose that
$s : P^\bullet \to (P')^\bullet$ is a quasi-isomorphism of bounded
above complexes consisting of objects of $\mathcal{P}$. Then
$F(P^\bullet) \to F((P')^\bullet)$ is a quasi-isomorphism because
$F(C(s)^\bullet)$ is acyclic by assumption (2). This already shows that
$LF$ is defined on $D^{-}(\mathcal{A})$ and that a bounded above
complex consisting of objects of $\mathcal{P}$ computes $LF$, see
Lemma \href{https://stacks.math.columbia.edu/tag/06XN}{lemma find existence computes (Tag 06XN)}.

\medskip\noindent
Next, let $K^\bullet$ be an arbitrary complex of $\mathcal{A}$.
Choose a diagram
$$
\xymatrix{
P_1^\bullet \ar[d] \ar[r] & P_2^\bullet \ar[d] \ar[r] & \ldots \\
\tau_{\leq 1}K^\bullet \ar[r] & \tau_{\leq 2}K^\bullet \ar[r] & \ldots
}
$$
as in Lemma \href{https://stacks.math.columbia.edu/tag/06XX}{lemma special direct system (Tag 06XX)}. Note that
the map $\colim P_n^\bullet \to K^\bullet$ is a quasi-isomorphism
because colimits over $\mathbf{N}$ in $\mathcal{A}$ are exact
and $H^i(P_n^\bullet) = H^i(K^\bullet)$ for $n > i$. We claim that
$$
F(\colim P_n^\bullet) = \colim F(P_n^\bullet)
$$
(termwise colimits) is $LF(K^\bullet)$, i.e., that $\colim P_n^\bullet$
computes $LF$. To see this, by Lemma \href{https://stacks.math.columbia.edu/tag/06XN}{lemma find existence computes (Tag 06XN)},
it suffices to prove the following claim. Suppose that
$$
\colim Q_n^\bullet = Q^\bullet
\xrightarrow{\ \alpha\ }
P^\bullet = \colim P_n^\bullet
$$
is a quasi-isomorphism of complexes, such that each
$P_n^\bullet$, $Q_n^\bullet$ is a bounded above complex whose terms are
in $\mathcal{P}$ and the maps $P_n^\bullet \to \tau_{\leq n}P^\bullet$ and
$Q_n^\bullet \to \tau_{\leq n}Q^\bullet$ are quasi-isomorphisms.
Claim: $F(\alpha)$ is a quasi-isomorphism.

\medskip\noindent
The problem is that we do not assume that $\alpha$ is given as a colimit
of maps between the complexes $P_n^\bullet$ and $Q_n^\bullet$. However,
for each $n$ we know that the solid arrows in the diagram
$$
\xymatrix{
& R^\bullet \ar@{..>}[d] \\
P_n^\bullet \ar[d] &
L^\bullet \ar@{..>}[l] \ar@{..>}[r] &
Q_n^\bullet \ar[d] \\
\tau_{\leq n}P^\bullet \ar[rr]^{\tau_{\leq n}\alpha} & &
\tau_{\leq n}Q^\bullet
}
$$
are quasi-isomorphisms. Because quasi-isomorphisms form a multiplicative
system in $K(\mathcal{A})$ (see Lemma \href{https://stacks.math.columbia.edu/tag/05RT}{lemma acyclic (Tag 05RT)})
we can find a quasi-isomorphism
$L^\bullet \to P_n^\bullet$ and map of complexes $L^\bullet \to Q_n^\bullet$
such that the diagram above commutes up to homotopy. Then
$\tau_{\leq n}L^\bullet \to L^\bullet$ is a quasi-isomorphism.
Hence (by the first part of the proof) we can find a bounded above
complex $R^\bullet$ whose terms are in $\mathcal{P}$ and a quasi-isomorphism
$R^\bullet \to L^\bullet$ (as indicated in the diagram). Using the result
of the first paragraph of the proof we see that
$F(R^\bullet) \to F(P_n^\bullet)$ and $F(R^\bullet) \to F(Q_n^\bullet)$
are quasi-isomorphisms. Thus we obtain an isomorphism
$H^i(F(P_n^\bullet)) \to H^i(F(Q_n^\bullet))$ fitting into the commutative
diagram
$$
\xymatrix{
H^i(F(P_n^\bullet)) \ar[r] \ar[d] &
H^i(F(Q_n^\bullet)) \ar[d] \\
H^i(F(P^\bullet)) \ar[r] &
H^i(F(Q^\bullet))
}
$$
The exact same argument shows that these maps are also compatible
as $n$ varies. Since by (4) and (5) we have
$$
H^i(F(P^\bullet)) =
H^i(F(\colim P_n^\bullet)) =
H^i(\colim F(P_n^\bullet)) = \colim H^i(F(P_n^\bullet))
$$
and similarly for $Q^\bullet$ we conclude that
$H^i(\alpha) : H^i(F(P^\bullet) \to H^i(F(Q^\bullet)$ is an isomorphism
and the claim follows.
\end{proof}
\end{document}
